% =====================================================================
%  CARNET D'ENTRAÎNEMENT — FICHE 3 : Trigonométrie
%  Thème 1 — Conversions d'angles & cercle trigonométrique
%  TUTORIEL
% =====================================================================
\documentclass[11pt,a4paper]{article}
\usepackage[utf8]{inputenc}
\usepackage[T1]{fontenc}
\usepackage{lmodern}
\usepackage[french]{babel}
\usepackage{amsmath,amssymb,mathtools}
\usepackage{icomma}
\usepackage{xcolor}
\usepackage{colortbl}
\usepackage{tikz}
\usepackage[most]{tcolorbox}
\usepackage{enumitem}
\usepackage{array}
\usepackage{booktabs}
\usepackage[a4paper,margin=2.2cm]{geometry}
\usetikzlibrary{arrows.meta,calc,angles,quotes,decorations.markings}

\definecolor{primary}{RGB}{214,51,108}
\definecolor{primarydark}{RGB}{140,20,64}
\definecolor{defcol}{RGB}{0,114,178}
\definecolor{excol}{RGB}{0,135,90}
\definecolor{methcol}{RGB}{90,90,100}
\definecolor{attncol}{RGB}{200,40,55}
\definecolor{sincol}{RGB}{205,40,55}
\definecolor{coscol}{RGB}{20,110,200}

\DeclareMathOperator{\tg}{tg}
\DeclareMathOperator{\cotg}{cotg}
\newcommand{\degr}{^{\circ}}
\newcommand{\Z}{\mathbb{Z}}

\tcbset{
  boxbase/.style={enhanced,breakable,boxrule=0.9pt,arc=2.5pt,
     left=3mm,right=3mm,top=2.3mm,bottom=2.3mm,
     fonttitle=\bfseries,coltitle=white,toptitle=0.6mm,bottomtitle=0.6mm},
}
\newtcolorbox{idee}[1][]{boxbase,colback=defcol!6,colframe=defcol,
  title={\textsc{L'essentiel}\ifx\relax#1\relax\relax\else\textmd{ : \itshape #1}\fi}}
\newtcolorbox{exemple}[1][]{boxbase,colback=excol!6,colframe=excol,
  title={\textsc{Exemple}\ifx\relax#1\relax\relax\else\textmd{ : \itshape #1}\fi}}
\newtcolorbox{methode}[1][]{boxbase,colback=methcol!8,colframe=methcol,sharp corners,
  title={\textsc{Méthode}\ifx\relax#1\relax\relax\else\textmd{ --- \itshape #1}\fi}}
\newtcolorbox{piege}[1][]{boxbase,colback=attncol!7,colframe=attncol,
  title={\textsc{Piège à éviter}\ifx\relax#1\relax\relax\else\textmd{ : \itshape #1}\fi}}

\newcommand{\heading}[1]{\medskip\noindent{\large\bfseries\color{primarydark}#1}\par\smallskip}

\pagestyle{empty}
\setlength{\parindent}{0pt}
\setlength{\parskip}{3pt}

\begin{document}

\begin{center}
{\LARGE\bfseries\color{primary}Thème 1 — Conversions \& cercle trigonométrique}\\[2pt]
{\color{primarydark}\textsc{Carnet d'entraînement — Fiche 3 : Trigonométrie \quad•\quad Tutoriel}}
\end{center}
\hrule height 1pt
\medskip

Ce thème installe les \textbf{réflexes de base} de la trigonométrie : convertir un angle, le placer
sur le cercle, lire son \textbf{quadrant}, le \textbf{signe} de ses nombres trigonométriques et ses
\textbf{valeurs remarquables}. Ces gestes reviennent dans presque tous les QCM : rendez-les
\emph{automatiques}.

\heading{1. Convertir degrés $\leftrightarrow$ radians}
Un demi-tour vaut $180\degr$ \emph{ou} $\pi$ radians. Tout part de cette égalité, résumée par la
règle de trois :
\[
\boxed{\;\dfrac{\alpha_{\text{deg}}}{180}=\dfrac{\alpha_{\text{rad}}}{\pi}\;}
\qquad\Longrightarrow\qquad
\alpha_{\text{rad}}=\alpha_{\text{deg}}\times\dfrac{\pi}{180},
\qquad
\alpha_{\text{deg}}=\alpha_{\text{rad}}\times\dfrac{180}{\pi}.
\]

\begin{exemple}[le $\pi$ se simplifie toujours]
$30\degr = 30\times\dfrac{\pi}{180}=\dfrac{\pi}{6}$ \ ;\qquad
$\dfrac{3\pi}{4}=\dfrac{3\pi}{4}\times\dfrac{180}{\pi}=\dfrac{3\times180}{4}=135\degr$.
Si le $\pi$ ne se simplifie pas, c'est le signe d'une erreur.
\end{exemple}

Les conversions à connaître par cœur (elles reviennent partout) :

\begin{center}
\renewcommand{\arraystretch}{1.35}
\begin{tabular}{|c|c|c|c|c|c|c|c|c|}
\hline
\rowcolor{primary!12}
\textbf{deg} & $0\degr$ & $30\degr$ & $45\degr$ & $60\degr$ & $90\degr$ & $180\degr$ & $270\degr$ & $360\degr$\\
\hline
\textbf{rad} & $0$ & $\frac{\pi}{6}$ & $\frac{\pi}{4}$ & $\frac{\pi}{3}$ & $\frac{\pi}{2}$ & $\pi$ & $\frac{3\pi}{2}$ & $2\pi$\\
\hline
\end{tabular}
\end{center}

\heading{2. Le cercle, les quadrants et le sens trigonométrique}
Le \textbf{cercle trigonométrique} a pour centre $O$ et pour rayon $1$. On tourne dans le
\textbf{sens trigonométrique} (inverse des aiguilles d'une montre = sens positif). Les deux axes le
découpent en quatre \textbf{quadrants} :
\[
\text{I : } 0\to\tfrac{\pi}{2}\quad|\quad
\text{II : } \tfrac{\pi}{2}\to\pi\quad|\quad
\text{III : } \pi\to\tfrac{3\pi}{2}\quad|\quad
\text{IV : } \tfrac{3\pi}{2}\to2\pi.
\]

\begin{center}
\begin{tikzpicture}[scale=2.1,line cap=round,>=Stealth]
  \fill[primary!7] (0,0)--(1,0) arc(0:90:1)--cycle;
  \fill[primary!13] (0,0)--(0,1) arc(90:180:1)--cycle;
  \fill[primary!7] (0,0)--(-1,0) arc(180:270:1)--cycle;
  \fill[primary!13] (0,0)--(0,-1) arc(270:360:1)--cycle;
  \draw[primary,line width=1.1pt] (0,0) circle (1);
  \draw[->,black!75] (-1.3,0)--(1.45,0) node[right]{\scriptsize$x=\cos$};
  \draw[->,black!75] (0,-1.3)--(0,1.4) node[above]{\scriptsize$y=\sin$};
  \draw[->,thick,primarydark] (1.15,0.28) arc (14:150:1.18);
  \node[primarydark] at (1.1,0.82) {\scriptsize sens $+$};
  \node[primarydark,font=\bfseries] at (45:0.55) {I};
  \node[primarydark,font=\bfseries] at (135:0.55) {II};
  \node[primarydark,font=\bfseries] at (225:0.55) {III};
  \node[primarydark,font=\bfseries] at (315:0.55) {IV};
  \foreach \a/\d in {0/0\degr,90/90\degr,180/180\degr,270/270\degr}
     \node[black!70,font=\scriptsize] at (\a:1.2) {$\d$};
  \fill (0,0) circle (0.02);
\end{tikzpicture}
\end{center}

\begin{idee}[un point du cercle porte cos et sin]
Au point $P$ associé à l'angle $\alpha$ : $P=(\cos\alpha\,;\,\sin\alpha)$. Donc \textbf{le cosinus se
lit horizontalement} (abscisse) et \textbf{le sinus verticalement} (ordonnée). Comme le rayon vaut
$1$ : $-1\le\cos\alpha\le1$ et $-1\le\sin\alpha\le1$.
\end{idee}

\heading{3. Signe selon le quadrant}
On lit le signe directement sur le cercle : $\cos$ (horizontal) est positif \emph{à droite},
$\sin$ (vertical) positif \emph{en haut}. La tangente $\tg\alpha=\frac{\sin\alpha}{\cos\alpha}$ est
positive quand $\sin$ et $\cos$ ont le \emph{même} signe.

\begin{center}
\renewcommand{\arraystretch}{1.3}
\begin{tabular}{|c|c|c|c|c|}
\hline
\rowcolor{primary!15}
\textbf{Quadrant} & I & II & III & IV\\
\hline
$\cos$ & $+$ & $-$ & $-$ & $+$\\ \hline
$\sin$ & $+$ & $+$ & $-$ & $-$\\ \hline
$\tg,\ \cotg$ & $+$ & $-$ & $+$ & $-$\\
\hline
\end{tabular}
\end{center}

\heading{4. Valeurs remarquables du premier quadrant}
Astuce du sinus : écris $\dfrac{\sqrt0}{2},\dfrac{\sqrt1}{2},\dfrac{\sqrt2}{2},\dfrac{\sqrt3}{2},
\dfrac{\sqrt4}{2}$. Le cosinus est la \emph{même liste à l'envers}.

\begin{center}
\renewcommand{\arraystretch}{1.6}
\begin{tabular}{|>{\columncolor{primary!12}}c|c|c|c|c|c|}
\hline
\rowcolor{primary!20}
$\alpha$ & $0$ & $\frac{\pi}{6}\,(30\degr)$ & $\frac{\pi}{4}\,(45\degr)$ & $\frac{\pi}{3}\,(60\degr)$ & $\frac{\pi}{2}\,(90\degr)$\\
\hline
$\sin\alpha$ & $0$ & $\frac{1}{2}$ & $\frac{\sqrt2}{2}$ & $\frac{\sqrt3}{2}$ & $1$\\
\hline
$\cos\alpha$ & $1$ & $\frac{\sqrt3}{2}$ & $\frac{\sqrt2}{2}$ & $\frac{1}{2}$ & $0$\\
\hline
$\tg\alpha$ & $0$ & $\frac{\sqrt3}{3}$ & $1$ & $\sqrt3$ & $\nexists$\\
\hline
\end{tabular}
\end{center}

\begin{methode}[lire un nombre trigonométrique d'un angle quelconque]
\begin{enumerate}[leftmargin=1.6em,topsep=1pt,itemsep=1pt]
  \item Ramener l'angle dans $[0\degr;360\degr[$ en ajoutant/retranchant des tours ($\pm360\degr$).
  \item Repérer le \textbf{quadrant} $\Rightarrow$ en déduire le \textbf{signe}.
  \item Trouver l'angle de référence (angle aigu avec l'axe horizontal) et lire sa valeur au tableau.
  \item Recoller le signe de l'étape 2.
\end{enumerate}
\end{methode}

\begin{exemple}[$\cos120\degr$]
$120\degr$ est au quadrant II ($\cos<0$). Référence : $180\degr-120\degr=60\degr$, et
$\cos60\degr=\frac12$. Donc $\cos120\degr=-\frac12$.
\end{exemple}

\begin{piege}[erreurs classiques]
$\bullet$ Oublier de simplifier le $\pi$ dans une conversion.\quad
$\bullet$ Confondre les lignes $\sin$ et $\cos$ (le sinus \emph{monte} de $0$ à $1$, le cosinus
\emph{descend} de $1$ à $0$).\quad
$\bullet$ Donner un signe $+$ à un cosinus du quadrant II ou III.
\end{piege}

\end{document}
